\section{Bounded formulas and Booleanization}
\label{sec:booleanization}

Quantifier-free induction is easiest to use when all finite side conditions are
represented by equations between primitive-recursive terms.  Bounded quantifiers do
not literally belong to the quantifier-free fragment, but every bounded formula has a
primitive-recursive characteristic function.  This familiar fact must be developed
constructively and in a form compatible with the embedded syntax.

\subsection{Bounded formula grammar}

Fix free variables \(\vec x\).  The bounded formulas are generated by
\[
\begin{split}
 A ::= {}& s=t \mid \bot \mid A\wedge A \mid A\vee A \mid A\to A \\
         &\mid \forall u<b(\vec x)\,A(u,\vec x)
          \mid \exists u<b(\vec x)\,A(u,\vec x),
\end{split}
\]
where all terms denote primitive-recursive functions.  Negation is \(A\to\bot\).
Because equality of naturals is decidable in \(\iPRA\), every such formula is
decidable, but we need more: a fixed characteristic term and embedded equivalence
proofs.

\subsection{Boolean arithmetic}

Let \([s=t]\in\B\) be the primitive-recursive equality test.  The checked
\code{Core/Pairing} module already supplies an equality term
\code{eqb} with object proofs for equal and unequal arguments.  Let
\(\sg(n)=0\) when \(n=0\) and \(1\) otherwise, and let
\(1\mathbin{\dot-}n\) denote truncated subtraction.

For Boolean inputs \(a,b\in\B\), define
\[
\begin{array}{rclcrcl}
 a\mathbin{\&}b&=&ab, &&
 a\mathbin{|}b&=&\sg(a+b),\\
 a\Rightarrow b&=&\sg((1\mathbin{\dot-}a)+b),&&
 \neg a&=&1\mathbin{\dot-}a.
\end{array}
\]
All four results are Boolean.  Intuitionistically,
\(a\Rightarrow b=1\) is equivalent to \(a=1\to b=1\): if \(a=0\), the
truncated difference contributes \(1\); if \(a=1\), the expression has sign
\(b\).  This finite case split is derivable from the successor axioms and
quantifier-free induction, and no excluded middle for arbitrary formulas is involved.

\subsection{Structural characteristic compiler}

\begin{construction}[Characteristic term]
Associate to every bounded formula \(A\) a primitive-recursive term
\(\chi_A(\vec x)\) by structural recursion:
\[
\begin{array}{rcl}
\chi_{s=t}&=&[s=t],\\
\chi_{\bot}&=&0,\\
\chi_{A\wedge B}&=&\chi_A\chi_B,\\
\chi_{A\vee B}&=&\sg(\chi_A+\chi_B),\\
\chi_{A\to B}&=&\sg((1\mathbin{\dot-}\chi_A)+\chi_B),\\
\chi_{\exists u<b A(u)}&=&\sg\!\left(\displaystyle\sum_{u<b}\chi_A(u)\right),\\
\chi_{\forall u<b A(u)}&=&\displaystyle\prod_{u<b}\chi_A(u).
\end{array}
\]
Empty sums are \(0\) and empty products are \(1\).  Bounded sums and products are
primitive recursive by ordinary recursion on \(b\).
\end{construction}

\begin{theorem}[Bounded Booleanization]
\label{thm:booleanization}
For every bounded formula \(A(\vec x)\), \(\iPRA\) derives uniformly:
\begin{align}
  &\chi_A(\vec x)=0\ \vee\ \chi_A(\vec x)=1,\label{eq:boolean-range}\\
  &A(\vec x)\ \longrightarrow\ \chi_A(\vec x)=1,\label{eq:to-char}\\
  &\chi_A(\vec x)=1\ \longrightarrow\ A(\vec x).\label{eq:from-char}
\end{align}
Consequently \(A\) is interchangeable with the quantifier-free equation
\(\chi_A=1\) in any finite trace invariant.
\end{theorem}

\begin{proof}
Proceed by structural induction on \(A\).  Equality and falsity are immediate from the
equality test and the successor axiom \(0\ne1\).  For conjunction, the induction
hypotheses make both factors Boolean.  Their product is \(1\) exactly when both are
\(1\); if either is \(0\), the product is \(0\).  For disjunction, the sign of the sum
is \(1\) exactly when at least one Boolean summand is \(1\).  These arguments use only
finite case analysis on numerical Booleans.

For implication, split on the Boolean value of \(\chi_A\).  If it is \(0\), the
characteristic value is \(1\), and \(A\to B\) follows because \(A\) would imply
\(\chi_A=1\), contradicting \(0=1\).  If \(\chi_A=1\), the characteristic value equals
\(\chi_B\), so the desired equivalence follows from the induction hypotheses for
\(B\).  This is an intuitionistically valid proof because the split is on a decidable
numerical equality, not on the arbitrary formula \(A\).

For the existential case, use induction on the bound \(b\) to prove that the bounded
sum is zero exactly when all characteristic values below \(b\) are zero, and has
positive sign exactly when some value is one.  The forward direction chooses the least
index whose summand is nonzero by bounded search; because each summand is Boolean, it
is one.  The reverse direction is immediate because a one summand makes the sum
positive.  Apply the induction hypothesis for the body to move between
\(\chi_A(u)=1\) and \(A(u)\).

For the universal case, induction on \(b\) proves that the bounded product is one
exactly when every factor below \(b\) is one.  The successor equation is
\[
   \prod_{u<b+1}\chi_A(u)
   =\left(\prod_{u<b}\chi_A(u)\right)\chi_A(b).
\]
A product of two Booleans is one exactly when both factors are one.  Again apply the
body induction hypothesis.  All inductions just mentioned have quantifier-free
invariants after replacing their bounded assertions by the corresponding sum or
product equations.  Thus they are available in \(\iPRA\).
\end{proof}

\subsection{Substitution and compilation discipline}

For use in a proof compiler, \cref{thm:booleanization} must be stable under
substitution.  If \(\sigma\) maps free variables to object terms, structural
construction gives
\[
   \chi_{A[\sigma]}=\chi_A[\sigma]
\]
up to the defining equations of the generated primitive-recursive symbols.  The Lean
implementation should prove this simultaneously with Booleanization, because the
bounded quantifier cases shift de Bruijn indices and are otherwise a common source of
errors.

The proposed \code{BForm} datatype in
\code{lean/FRNF/ObjectTargets.lean} records the grammar, with fields
\code{toFm} and \code{char}.  The three intended endpoints are exactly
\cref{eq:boolean-range,eq:to-char,eq:from-char}.  Completing them should reuse the
checked bounded-sum infrastructure rather than construct a second arithmetic library.

\subsection{Why implication is included}

One could normalize every bounded formula to positive Boolean combinations of atomic
equalities, but keeping implication is useful.  Contracts frequently have the form
``if the index is in range, then the table entry is valid.''  Compiling implication
directly avoids a separate negation-normal-form pass.  The construction remains
constructive because the antecedent has already been compiled to a Boolean.

\subsection{Consequence for induction}

Let \(I(t,\vec x)\) be a bounded assertion about the state of a loop at time \(t\).
Replace it by
\[
   \chi_I(t,\vec x)=1.
\]
The latter is a literal equation between primitive-recursive terms and hence a
quantifier-free formula.  If object derivations establish the base and successor
steps, the \(\PRAAx\) quantifier-free induction axiom closes the loop.  This is the
basic bridge from finite mathematics to the restricted induction of \(\iPRA\).
